\chapter{The complete positive-total-scale canonical source union}
\label{component:pc0b2849f8d09}
\noindent\textbf{Component }\path{pc0b2849f8d09}\par
{\small\noindent Source: \path{NEW/global/positive_scale_union_REDERIVED.tex}\par}
{\small\noindent SHA256: \path{c0b2849f8d092c22f3c525128bb352681aa7e80818bfcf91489a6521c1bfc33c}\par}
{\small\noindent Inventory use: AU1-AU19 complete positive-total-scale canonical source union and proof-use scope strengthening\par}
\medskip
\section{The complete positive-total-scale canonical source union}
\label{pc0b2849f8d09:sec:positive-scale-union-rederived}

\paragraph{Exact statement.}
Let \(K_c,K_t\) be labelled real closed extensions of \(\mathbb R\).
The original canonical coefficient vector is physical over \(K_c\),
has nonnegative coefficients, and has an actual finite noncommuting
zero contact. In C1 the exact omissions are \(\lambda=n=0\),
with \(m,C,A,\alpha>0\), \(D,\beta\geq0\), \(D+\beta>0\);
in C2 they are \(\lambda=\beta=0\), with
\(m,C,A,D,\alpha>0\), \(n\geq0\).
Use the actual native sign chart of this contact, with its explicitly
positive source gaps, as supplied by canonical entry.
Assume that normalization by \(d>0\) makes all coefficients bounded.
Equivalently the normalized total-scale coordinate
\(d/(d+\lambda+n+m+C+A+D+\alpha+\beta)\) is a positive unit.
Then the entire original homogeneous slack is a sum of squares
in \(K_c\otimes_{\mathbb R}K_t\) for every independent projective
testing pair and every real or complex testing state.
No prior-admission complement is assumed.
This is the A5 terminal union at the explicit canonical-entry interface;
global entry into that interface is a separate theorem.

Division by \(d\) is a positive pure-source normalization.
Hence take \(d=1\) below. Every nonnegative bounded source quantity
is either infinitesimal (including literal zero where permitted)
or a positive unit. The proof makes that finite split exhaustive,
rather than using undefined historical terminal labels.

\subsection{An explicit noncommuting-residue lemma for this union}

For \(\lambda=0\) the determinant Laurent polynomial is
\[
 L(U,V)=h+\frac{q_{01}}U+\frac{q_{02}V}U+\frac{q_{10}}V
       +q_{12}V+\frac{q_{20}U}V+q_{21}U+q_{22}UV ,
 \qquad U=u^2,\ V=v^2.
 \tag{AU1}
\]
All noncentral coefficients are nonnegative, with the exact
factorization in the physical-core input.
Suppose \(C,A,D,\alpha\) are units.
If \(m\) is weak, then \(\bar q_{01}=\bar q_{02}=0\), whereas
\(\bar q_{21}>0\). At a finite nonzero real \(U,V\),
\[
 U\partial_U\bar L
   =\bar q_{20}U/V+\bar q_{21}U+\bar q_{22}UV>0.
 \tag{AU2}
\]
Thus a physical residue source has no zero there:
a differentiable nonnegative function at an interior zero
has zero derivative.

If \(m\) is a unit, then \(q_{01},q_{10},q_{12},q_{21}\) are
positive units. At any real noncommuting zero, the Hessian of
\(L(e^s,e^t)\) is
\[
 \sum_{\gamma} \bar q_\gamma
             e^{\gamma_1s+\gamma_2t}\gamma\gamma^{\mathsf T}
 \succ0 .
 \tag{AU3}
\]
The terms with exponents \((-1,0),(0,-1)\) already span both
directions. This is a real residue calculation only.

Here is the full-parent consequence, without a mixed inverse.
The \(Z=0\) principal state block is positive definite:
\(A Y^2\), \(C(uW-Y)^2\), and the positive core \(X\)-square
already force \(W=X=Y=0\), since \(u\ne0\).
A zero state therefore has \(Z\ne0\), and it can be normalized
to \(Z=1\) in the testing field. Its real three-state Schur
complement is
\[
 g(u,v)=\frac{\det P(u,v)}
                 {\det P(u,v)|_{Z=0}}
 =\frac{D_u^2D_v^2UV}{16\det P(u,v)|_{Z=0}}\,L(U,V)
 \tag{AU4}
\]
for the quarter-parent matrix \(P\).
At a zero the value and first derivatives of \(L\) vanish.
AU3 and the nonzero real change from \(u,v\) to \(\log U,\log V\)
therefore give a positive angular Hessian for \(g\).
The full Hessian in the three state coordinates and two angles
is positive definite.

The polynomial implicit-function theorem supplies, in the source
field only, a critical center near this real five-variable zero.
One may justify this over an arbitrary real closed source extension
by the semialgebraic local implicit-function theorem over
\(\mathbb R\) and real closed transfer. The Hessian is a real
positive definite matrix at the base point, so its invertibility
and the local uniqueness condition are explicit.
The original source is physical, hence its value \(\mu\) at
that actual source-field center is nonnegative.
The exact finite Taylor expansion is
\[
 F=\mu+N^{\mathsf T}HN,\qquad
 \bar H=\tfrac12\operatorname{Hess}F_0\succ0 .
 \tag{AU5}
\]
The five normals are differences of source coordinates and
independent testing coordinates with the same real residues.
All coefficients of \(H\) are bounded. The separated strict-Gram
lemma proves AU5 tensor-SOS; the pure-source \(\mu\) is a square.
Positive residue slack is the ordinary scalar strict-Gram exit.
This proves every noncommuting-residue test on this strong source
domain. It does not import a broad noncommuting theorem by name.

One further residue case is needed.
In C2 let \(\alpha\) be weak and let each of \(n,m,C,A,D\) be a unit.
Then
the residue Laurent polynomial is
\[
 \bar L=h+q_{02}V/U+q_{20}U/V+q_{12}V+q_{21}U .
\]
Its radial derivative is
\[
 (U\partial_U+V\partial_V)\bar L=q_{12}V+q_{21}U>0.
 \tag{AU6}
\]
It has no interior zero. Thus its finite noncommuting state
matrices are positive definite. This explicitly supplies the
noncommuting input in the retained finite-ratio proof as well.

\subsection{Two retained full-parent inputs, with all hypotheses fitted}

\paragraph{Strong and finite-edge sources.}
The actual retained proof \emph{finite-pitchfork body} establishes
all projective tests when \(d,C,A,D,\alpha\) are units, with
arbitrary nonnegative \(n,m,\beta\) and \(\lambda=0\).
Its noncommuting call is replaced here by AU2--AU5.
The only non-strict finite commuting edge has \(u=0\), \(v_0\ne0\),
weak \(m\), and the \(W\) state ray. On that residue face,
\[
 Q(U,V)=U E(V)+U^2B(V),\quad
 E(V)=q_{10}+hV+q_{12}V^2,\quad
 B(V)=q_{20}+q_{21}V+q_{22}V^2,
 \tag{AU7}
\]
where \(q_{10},q_{12},q_{21}>0\).
Physicality makes \(E(V)\geq0\) on \(V>0\).
Thus either \(E(v_0^2)>0\), giving a strict four-normal Gram,
or \(E(V)=q_{12}(V-v_0^2)^2\).
The displayed three-state complement in the actual
\emph{finite-edge terminal} is positive: its lower determinant
has the positive factor
\[
 \alpha+2\alpha C+2\alpha A+2\alpha D(1+v_0^2)
                     +Dv_0^2(1+2C+2A)>0.
\]
Also \(B(v_0^2)>0\). Every hypothesis of that retained local
theorem is therefore verified on the original source.

For specificity, its pure-testing event normalization makes
\(vW-X+uvY-uZ=1\). The complete parent becomes
\(4m+u^2K(x,a,b,v,u^2)\) only as a formal source-center calculation.
The actual tensor normals are
\[
 X-u^2x_c,\quad Y-ua_c,\quad Z-ub_c,\quad
 u(v-v_c),\quad u^2-U_c .
 \tag{AU8}
\]
The divided \(x\)-equation and the full three-variable Hessian
are strict; the last scalar pivot is the positive
\(16(1+v_0^2)^2B(v_0^2)/(v_0^2\det M_0)\).
The two signs of the optimized source constant give either
the nonnegative source value at \(U_c>0\), or the retained
\(4m+\kappa u^2\) at \(U_c=0\).
The retained theorem expands all first state normals and the
complete state-quadratic tail. Its finite monic divisions do
not divide a testing infinitesimal.

The same actual body supplies the four literal projective
corners. In quarter-parent normalization their residue forms
after the indicated whole event is retained are
\[
\begin{array}{c|l}
 (0,0),\ W=1,\ m\text{ retained}
 &(n+\beta+\frac12)X^2+K(u-Y)^2+AY^2+\alpha Z^2\\
 (\infty,0),\ X=1,\ n\text{ retained}
 &CW^2+AY^2+(D+\frac12)(W+v)^2+
     \alpha Z^2+(m+\beta+\frac12)(s+Z)^2\\
 (0,\infty),\ W=1,\ m\text{ retained}
 &(n+\beta+\frac12)X^2+C(u-Y)^2+
     (A+\alpha+\frac12)Y^2+DZ^2\\
 (\infty,\infty),\ Z=1,\ \beta\text{ retained}
 &CW^2+AY^2+D(X-s)^2+
     (n+\frac12)X^2+(m+\alpha+\frac12)(Y-b)^2 .
\end{array}\tag{AU9}
\]
Here \(s=1/u,b=1/v\) in the reciprocal charts,
\(K=C+D+\frac12\), and coefficient letters denote residues.
Each displayed quadratic form is strict on its printed normal
variables. The entire remainders belong to their square ideals,
as shown by the actual polynomial expansions in the retained body.
The two single-infinity edges have, among their positive terms,
\(\frac12[(W+vX)^2+(vY-Z)^2]+C D_vW^2+A D_vY^2\)
or \(\frac12D_u(X^2+Y^2)+C(uW-Y)^2+D(uX-Z)^2\),
respectively, and are strict when the other angle is nonzero.
These forms and AU7--AU8 exhaust all states and angles.
We denote this exact all-tests input by \(\mathrm{S}\).

\paragraph{The unit-\(m,n\) weak-alpha source.}
The actual retained \emph{weak-alpha unit-\(n\) body} assumes
\(d,n,m,C,A,D\) all positive units and \(\alpha\) weak in C2.
AU6 supplies its noncommuting input.
Its only remaining non-strict angle corner is \(u=v=0\).
At \(W=1,Z\) infinitesimal, retain the whole \(\alpha\) packet;
the quadratic form
\[
 (n+\tfrac12)X^2+m(v-X)^2+K(u-Y)^2+AY^2
 \tag{AU10}
\]
is strict in \(X,Y,u,v\).
At \(Z=1,W\) bounded the full quadratic jet \(Q_w\) in the
retained body is positive semidefinite by physicality.
Its \(X,Y\) block is positive, and it has no nonzero kernel
with either \(u=0\) or \(v=0\): the positive \(n,m,A,D\) terms
force all remaining variables to vanish in each case.
Hence a singular angular Schur complement has a one-dimensional
kernel with a finite nonzero real ratio \(t_0=u/v\).

The full homogeneous state jet at that ratio has a positive
\(W,X,Y\) principal complement, already from
\(nX^2+m(W-X)^2+AY^2\).
Moreover
\[
 q_{20},q_{02}>0,\quad
 q_{20}t_0^4+h t_0^2+q_{02}=0,\quad
 q_{21}t_0^4+q_{12}t_0^2>0.
 \tag{AU11}
\]
These are exactly the actual \emph{finite-ratio corner terminal}
hypotheses. Its source-only critical centers and complete
normal-tail allocation use the undivided tensor rows
\[
 X-vx_c,\quad Y-vy_c,\quad u-vt_c,\quad
 v(W-w_c),\quad v^2-V_c .
\]
The source-selected \(t_0\) imposes no restriction on the actual
independent quotient \(u/v\). The explicit 55-monomial allocation
retains \(W-w_c\) in bounded coefficients rather than replacing it
by a mixed quotient. Its strict four-variable Hessian and positive
last pivot prove the whole parent.
The retained complementary-state argument and full projective
forms cover every other state, sign and literal endpoint.
We denote this exact all-tests input by \(\mathrm{CU}\).
Only the printed unit hypotheses are used; no weak-\(m\) or
weak-\(n\) conclusion is inferred from it.

\subsection{Exact symmetries and the source partition}

Both transformations below preserve the core and nonnegative event
coefficients, actual physicality, finite nonzero contact, and pure
total-scale normalization. They act separately in source and testing
fields, so transport a full tensor certificate.

For C1 let
\[
 \sigma_1=(\lambda\ n)(m\ C)(A\ \alpha)(D\ \beta),\qquad
 G_v(W,X,Y,Z)=(vW-X,-W-vX,Z-vY,Y+vZ).
\]
The entire raw parent has the polynomial identity
\[
 u^2F_{\sigma_1c}(1/u,v;G_v\psi)
       =(1+v^2)F_c(u,v;\psi),\qquad
 G_v^{\mathsf T}G_v=(1+v^2)I .
 \tag{AU12}
\]
Homogenization interchanges the two Alice angle coordinates;
the Bob-homogeneous \(G\) has norm factor \(v_0^2+v_1^2>0\).
Thus literal endpoints are transported without dividing zero.

For C2 let
\[
 \sigma_2=(\lambda\ \beta)(m\ D)(C\ \alpha),\qquad
 J(W,X,Y,Z)=(Z,X,Y,W).
\]
The full identity is
\[
 u^2v^2F_{\sigma_2c}(1/v,1/u;J\psi)=F_c(u,v;\psi).
 \tag{AU13}
\]
It is the physical party interchange with both first observables
negated, followed by the local real \(X\) basis changes.
The homogeneous identity includes both projective endpoints.
In particular the source squared half-angles transform as
\((U,V)\mapsto(1/V,1/U)\).

The two actual forbidden weak faces are
\[
\begin{aligned}
 h|_{\lambda=\beta=A=0}
 &=-(1+2n)(2D+2C+1)(2\alpha+1+2m),\\
 h|_{\lambda=n=m=C=0}
 &=-(2D+1)(2\beta+1)(2A+2\alpha+1).
\end{aligned}\tag{AU14}
\]
In C2 with \(A\) weak, the only positive Laurent terms are
\(q_{01}/U+q_{02}V/U+q_{12}V\), which tend to zero along
\(U=t^2,V=1/t,\ t\to\infty\).
In C1 with both \(m,C\) weak the only positive term is \(q_{10}/V\),
which tends to zero as \(V\to\infty\).
The negative constants contradict physicality.
Thus C2 has \(A\) unit, and C1 has at least one of \(m,C\) unit.
These are statements about physical residue sources, not an
assumption inherited from a prior routing complement.

\paragraph{C2.}
If \(\alpha\) is weak, WA1--WA6 force \(C,A,D\) to be units.
There are precisely three disjoint alternatives:
\[
\begin{array}{c|c}
 n,m\text{ units}&\mathrm{CU}\\
 n\text{ unit},\ m\text{ weak}&\mathrm{WA1\!-\!WA27}\\
 n\text{ weak},\ m\text{ arbitrary}&\mathrm{WN1\!-\!WN44}.
\end{array}\tag{AU15}
\]
If \(C\) is weak, apply \(\sigma_2\) and the same table.
Every unit requirement in the transformed table is a consequence
of the transformed physical source, including its required
\(D'=m\) bound; it is not assumed for an excluded original case.

It remains that \(C,A,\alpha\) are units.
The actual dictionary then has
\[
 q_{01}=ma_1,\ q_{02}=mDa_2,\ q_{10}=a_3,\
 q_{12}=Da_4,\ q_{20}=na_5,\ q_{21}=a_6 ,
\]
where every \(a_i\) is a positive bounded unit.
The contact equation bounds \(U,V^{-1},m/U,DV\).
The two angular equations give exactly
\[
 \begin{split}
 DV^2(a_4+a_2m/U)&=a_3+na_5U,\\
 U^2(a_6+na_5/V)&=m(a_1+a_2DV).
 \end{split}
\]
All parentheses are positive bounded units, so
\(DV^2\asymp1\), \(U^2\asymp m\).
Consequently the exhaustive all-tests dispatch is
\[
\begin{array}{c|c}
 D\text{ unit},\ m\text{ arbitrary}&\mathrm S\\
 D\text{ weak},\ m\text{ unit}&\sigma_2(\mathrm S)\\
 D,m\text{ weak}&\mathrm{RC1\!-\!RC25}.
\end{array}\tag{AU16}
\]
RC is the new direct full-parent C2 reciprocal theorem; it permits
arbitrary \(n\geq0\). No unit-\(n\) condition has entered this
part of the partition.

\paragraph{C1.}
If \(\alpha\) is weak, WN applies directly, with all its common
source-entry implications proved in WA1--WA10.
If \(A\) is weak, apply \(\sigma_1\), making \(\alpha'=A\) weak,
and use WN on the transformed physical source.
Otherwise \(A,\alpha\) are units.
By AU14, orient with \(\sigma_1\), if necessary, so that \(C\)
is a unit. The exact dictionary becomes
\[
\begin{gathered}
 q_{01}=ma_1,\ q_{02}=mDa_2,\ q_{10}=a_3,\quad
 q_{12}=\beta ma_4+Da_5,\quad q_{21}=a_6,\quad q_{22}=\beta a_7,\\
 U^2(a_6+\beta a_7V)=m(a_1+a_2DV),\\
 a_3=V^2\{mDa_2/U+\beta ma_4+Da_5+\beta a_7U\}.
\end{gathered}\tag{AU17}
\]
The contact equation again bounds \(U,V^{-1},m/U,DV\).
If \(m\) is weak, \(U\) is infinitesimal.
If \(m\) is a unit, \(U\) is a positive unit.
Indeed, the contact equation also bounds \(\beta UV\); if
\(U\) were infinitesimal, then
\(U^2(a_6+\beta a_7V)\) would be infinitesimal, contradicting
the positive unit \(ma_1\) on the right of AU17.
For unit \(m\), the last braces are comparable with \(D+\beta\).
For weak \(m\), bounded \(V\) forces \(m/U\) weak in the
preceding equation, and then the last equation forces \(D\)
unit; conversely unit \(D\) forces bounded \(V\).
Thus the following source dispatch is exhaustive:
\[
\begin{array}{c|c}
 D\text{ unit},\ m\text{ arbitrary}&\mathrm S\\
 D\text{ weak},\ m\text{ unit},\ \beta\text{ unit}
       &\sigma_1(\mathrm S)\\
 m\text{ unit},\ D,\beta\text{ weak}
       &\mathrm{PP1\!-\!PP19}\\
 m,D\text{ weak},\ \beta\text{ arbitrary}
       &\mathrm{CR1\!-\!CR39}.
\end{array}\tag{AU18}
\]
In the second line all of \(m,C,A,\alpha\) are units before
reflection, so the reflected strong theorem has every required
unit. PP is the new C1 projective-pitchfork theorem, with actual
contact \(U\) unit and \(V\) unbounded.
CR is the new C1 reciprocal all-tests theorem, including unit
\(\beta\) and literal allowed \(D=0\) or \(\beta=0\).
Its \(m,D\) weak scope is exactly the last line.
No boundary is promoted by a density argument.

AU15--AU18 cover every permitted bounded positive-\(d\) canonical
source. Each leaf proves the entire original parent at every
independent projective test, with the exact same source coefficients.
The full-parent symmetries, positive pure-source \(d\) scaling,
positive pure-testing homogeneous factors and complex-state split
restore the assertion in the opening paragraph. This completes
A5 at its canonical-entry interface.

\paragraph{Proof-use scope lemma.}
The bare source domains in this union do not silently discard
the prior-complement phrases retained in the earlier statements.
Their removal here follows by checking the actual proofs:
RC1 uses only the two equations in AU16; RC2--RC5 use the actual
kernel equations, positive native gaps, and the listed unit bounds;
RC6--RC25 are full-parent identities with those bounds.
PP1--PP5 use only AU17, the actual kernel and angle jets, and its
listed unit bounds; PP6--PP19 and its subsequent outer identities
use no source exclusion.
CR1--CR8 likewise derive every shape estimate from AU17 and the
actual source equations, and CR9--CR39 use only those estimates.
WA1--WA10 prove their own source entry, while WA11--WA27 and
WN1--WN44 require only the printed weak/unit alternatives.
The retained S and CU bodies state physical coefficient domains
directly and have no prior-admission premise.
Consequently the same proofs establish the implications
\[
 \text{physical native source and listed scale bounds}
 \ \Longrightarrow\ \text{all-tests certificate}
 \tag{AU19}
\]
without any earlier-terminal exclusion.
This is a proof-use strengthening on fixed native source charts;
it does not identify an arbitrary contact with such a chart.
That sign/native entry implication remains an explicit global
dependency.

\paragraph{Evidence and claim boundary.}
The dependency receipt identifies the actual finite-edge,
finite-ratio, fit and strict-Gram proof bodies, not just their
historical labels. The focused checker verifies both complete
parent symmetries, the forbidden faces, both exact stationarity
dictionaries and the finite weak/unit dispatch.
The source comparisons, full Hessian proof and typed all-tests
applications are the analytical argument above.
This theorem does not assert weak-total-scale C1/C2 closure,
global support-cone entry, or a final finite-dimensional lift.
Those require their separately proved source and testing unions.
